\documentclass[11pt,a4paper]{article}

% ------------------------------------------------------------------
% Paquetes (todos estándar, compatibles con Overleaf sin configuración extra)
% ------------------------------------------------------------------
\usepackage[utf8]{inputenc}
\usepackage[T1]{fontenc}
\usepackage[spanish]{babel}
\usepackage[margin=2.5cm]{geometry}
\usepackage{xcolor}
\usepackage{booktabs}
\usepackage{longtable}
\usepackage{array}
\usepackage{ragged2e}
\usepackage{enumitem}
\usepackage{hyperref}
\usepackage{fancyhdr}
\usepackage{titlesec}
\usepackage{colortbl}
\usepackage{multicol}

\definecolor{critico}{RGB}{178,34,34}
\definecolor{advertencia}{RGB}{184,115,10}
\definecolor{ok}{RGB}{30,120,90}
\definecolor{gris}{RGB}{90,90,90}
\definecolor{fondoheader}{RGB}{240,240,240}

\hypersetup{
  colorlinks=true,
  linkcolor=teal,
  citecolor=teal,
  urlcolor=teal,
  pdftitle={Gamma Lab - Informe de la suite de pruebas},
  pdfauthor={Auditoria de codigo}
}

\newcommand{\okmark}{\textcolor{ok}{\textbf{OK}}}
\newcommand{\errmark}{\textcolor{critico}{\textbf{ERROR}}}
\newcommand{\warnmark}{\textcolor{advertencia}{\textbf{ADVERTENCIA}}}

\newcolumntype{L}[1]{>{\small\RaggedRight\arraybackslash}p{#1}}

\pagestyle{fancy}
\fancyhf{}
\fancyhead[L]{Gamma Lab -- Informe de la suite de pruebas}
\fancyhead[R]{\thepage}
\renewcommand{\headrulewidth}{0.4pt}

\titleformat{\section}{\normalfont\Large\bfseries}{\thesection}{1em}{}
\titleformat{\subsection}{\normalfont\large\bfseries}{\thesubsection}{1em}{}

\title{\textbf{Gamma Lab}\\[4pt]\large Informe técnico: análisis de la suite de pruebas\\
\large Qué prueba cada test, cómo funciona y qué está mal}
\author{Auditoría de código}
\date{14 de agosto de 2026}

\begin{document}

\maketitle
\thispagestyle{empty}

\begin{abstract}
\noindent
Este informe documenta, archivo por archivo, la suite de pruebas automatizadas de \textbf{Gamma Lab}
(\texttt{test/}): qué verifica cada test, cómo está construido paso a paso, y qué errores concretos se
encontraron al ejecutarla realmente contra el código actual. La suite tiene tres carpetas:
\texttt{services\_\allowbreak test} (pruebas unitarias de servicios), \texttt{plugins\_\allowbreak test} (pruebas de integración
contra una referencia de MATLAB) y \texttt{performance\_\allowbreak test} (benchmarks de rendimiento, agregados durante
esta auditoría). Ningún archivo de \texttt{core/} o \texttt{plugins/} fue modificado para producir este
informe; todas las cifras provienen de ejecuciones reales de \texttt{pytest} sobre el repositorio.
\end{abstract}

\tableofcontents
\newpage

% ==================================================================
\section{Resumen ejecutivo}
% ==================================================================

Al ejecutar \texttt{pytest} sobre el repositorio tal como está hoy, el resultado es:

\begin{center}
\begin{tabular}{lccc}
\toprule
\textbf{Carpeta} & \textbf{Archivos} & \textbf{Tests} & \textbf{Estado} \\
\midrule
\texttt{test/services\_\allowbreak test/} & 6 & 60 & \okmark\ (60/60 pasan) \\
\texttt{test/plugins\_\allowbreak test/}  & 8 & 24 & \errmark\ (0/24 se ejecutan -- error de importación) \\
\texttt{test/performance\_\allowbreak test/} & 2 & 10 & \okmark\ (10/10 pasan) \\
\bottomrule
\end{tabular}
\end{center}

\noindent\textbf{Total real medido:} \texttt{70 passed, 8 errors} (los 8 errores son archivos completos que
no logran ni importarse, no fallas de aserciones).

El hallazgo más importante: \textbf{los 8 archivos que validan los plugins de análisis contra una referencia
de MATLAB no se ejecutan en absoluto}, por un import roto que sobrevivió a un refactor de módulos. Al
corregir ese import en una copia aislada (fuera del repositorio, solo para diagnóstico -- el repositorio real
nunca fue modificado), 19 de 24 tests pasan y \textbf{5 fallan de verdad}, revelando errores numéricos reales
en los plugins \texttt{wavelet}, \texttt{wavelet\_\allowbreak average} y \texttt{psd\_\allowbreak average}.

% ==================================================================
\section{Arquitectura general de la suite}
% ==================================================================

\texttt{test/conftest.py} es el único archivo compartido por toda la suite. Hace dos cosas:

\begin{enumerate}
  \item Agrega la raíz del repositorio a \texttt{sys.path}, para que \texttt{import core...} y
        \texttt{import plugins...} funcionen sin instalar el proyecto como paquete.
  \item Registra el marcador \texttt{perf} de \texttt{pytest} (agregado en esta auditoría junto con la
        carpeta \texttt{performance\_\allowbreak test}, ver sección~\ref{sec:performance}) para que no aparezca un
        \textit{warning} de marcador desconocido.
\end{enumerate}

No hay más fixtures compartidas: cada archivo de test define sus propios \textit{helpers} y datos de
prueba, lo cual produce bastante duplicación de código entre archivos (se detalla por sección).

\subsection{Las tres capas}

\begin{itemize}
  \item \textbf{\texttt{services\_\allowbreak test/}} -- pruebas unitarias clásicas. Todo lo externo (PyQt5, VTK,
        \texttt{pyabf}, \texttt{pyedflib}) se sustituye por objetos \textit{Dummy} escritos a mano o por
        \texttt{unittest.mock.Mock}. No dependen de archivos reales ni de una interfaz gráfica.
  \item \textbf{\texttt{plugins\_\allowbreak test/}} -- pruebas de integración. Cargan el archivo ABF real de
        \texttt{test/data/17308005.abf}, cortan los ensayos (\textit{trials}) con la lógica real de
        \texttt{core/filters/trials.py}, ejecutan el método de cálculo real de cada plugin (usando una
        subclase ``sin cabeza'' que solo desactiva el renderizado VTK/Qt) y comparan el resultado contra un
        CSV de referencia generado en MATLAB.
  \item \textbf{\texttt{performance\_\allowbreak test/}} -- no verifica corrección, mide tiempo de ejecución del código
        real (cómputo y renderizado) para poder comparar ``antes'' y ``después'' de futuras optimizaciones.
\end{itemize}

\newpage
% ==================================================================
\section{\texttt{test/services\_\allowbreak test/} -- pruebas unitarias de servicios}
% ==================================================================

Los 6 archivos de esta carpeta suman 60 tests. Todos pasan, todos corren en aproximadamente 1 segundo en
total, y ninguno toca disco salvo para escribir archivos temporales de exportación (CSV/JSON) que luego se
leen y verifican.

% ------------------------------------------------------------------
\subsection{\texttt{test\_\allowbreak dataset.py} -- 7 tests}
% ------------------------------------------------------------------

\textbf{Qué prueba:} la clase \texttt{DataStore} (\texttt{core/services/data\_\allowbreak store.py}), el diccionario
central donde se guardan las señales cargadas y cuál está ``activa''.

\textbf{Cómo funciona, paso a paso:}
\begin{enumerate}
  \item Un \textit{helper} \texttt{make\_\allowbreak signal\_\allowbreak dataset(...)} construye un \texttt{SignalDataset} real
        (no un mock) con señales aleatorias de NumPy, para no depender de ningún archivo.
  \item Cada test crea un \texttt{DataStore()} nuevo (vacío) y ejercita una operación puntual.
  \item No hay mocks: es todo objeto real, así que si \texttt{DataStore} cambia su comportamiento interno,
        el test lo detecta de verdad.
\end{enumerate}

\begin{center}
\begin{longtable}{L{5.5cm}L{9cm}}
\toprule
\textbf{Test} & \textbf{Qué verifica} \\
\midrule
\endhead
\texttt{test\_\allowbreak set\_\allowbreak get\_\allowbreak has\_\allowbreak remove\_\allowbreak and\_\allowbreak items} & \texttt{set/get/has/remove/items} básicos del diccionario. \\
\texttt{test\_\allowbreak add\_\allowbreak signal\_\allowbreak with\_\allowbreak explicit\_\allowbreak key\_\allowbreak and\_\allowbreak get\_\allowbreak signals\_\allowbreak filters} & Al agregar una señal con clave explícita, \texttt{get\_\allowbreak signals()} la devuelve y excluye entradas no-señal. \\
\texttt{test\_\allowbreak add\_\allowbreak signal\_\allowbreak auto\_\allowbreak keys\_\allowbreak increment} & Sin clave explícita, las claves autogeneradas siguen el patrón \texttt{raw\_\allowbreak signal\_\allowbreak 1, \_\allowbreak 2, \_\allowbreak 3...} \\
\texttt{test\_\allowbreak set\_\allowbreak active\_\allowbreak signal\_\allowbreak success\_\allowbreak and\_\allowbreak getters} & Marcar una señal como activa y leerla de vuelta con los distintos getters. \\
\texttt{test\_\allowbreak set\_\allowbreak active\_\allowbreak signal\_\allowbreak invalid\_\allowbreak key\_\allowbreak raises} & Marcar como activa una clave inexistente lanza \texttt{ValueError}. \\
\texttt{test\_\allowbreak set\_\allowbreak active\_\allowbreak signal\_\allowbreak non\_\allowbreak signal\_\allowbreak raises} & Marcar como activo algo que no es un \texttt{SignalDataset} lanza \texttt{ValueError}. \\
\texttt{test\_\allowbreak clear\_\allowbreak active\_\allowbreak signal\_\allowbreak and\_\allowbreak is\_\allowbreak active\_\allowbreak signal} & \texttt{clear\_\allowbreak active\_\allowbreak signal()} deja el store sin señal activa. \\
\bottomrule
\end{longtable}
\end{center}

\textbf{Errores encontrados:} ninguno. Es, junto con \texttt{test\_\allowbreak trials.py}, el archivo mejor construido de
toda la suite: sin mocks innecesarios, aserciones concretas, nombres de test que describen exactamente el
caso cubierto.

% ------------------------------------------------------------------
\subsection{\texttt{test\_\allowbreak export.py} -- 8 tests}
% ------------------------------------------------------------------

\textbf{Qué prueba:} \texttt{ExportService} (\texttt{core/services/export\_\allowbreak service.py}), que exporta el
gráfico activo a imagen (PNG/JPEG/BMP/TIFF) o la tabla de datos a CSV/JSON/XLSX.

\textbf{Cómo funciona, paso a paso:}
\begin{enumerate}
  \item Un \textit{helper} \texttt{make\_\allowbreak service(tmp\_\allowbreak path, ...)} construye un \texttt{ExportService} real,
        pero con \texttt{parent}, \texttt{vtk\_\allowbreak widget} y \texttt{alerts} reemplazados por
        \texttt{unittest.mock.Mock()}.
  \item Para los tests de imagen, se usa \texttt{monkeypatch} para sustituir las clases de VTK
        (\texttt{vtkWindowToImageFilter}, \texttt{vtkPNGWriter}, etc.) por clases \texttt{Dummy*} que solo
        registran si fueron llamadas, sin generar un archivo real.
  \item Para los tests de tabla, se construyen objetos \texttt{DummyTable}/\texttt{DummyChart}/\texttt{DummyPlot}
        que imitan la interfaz mínima de VTK que \texttt{export\_\allowbreak table} necesita, y luego se \textbf{lee el
        archivo CSV/JSON real que quedó escrito en disco} y se compara su contenido.
\end{enumerate}

\begin{center}
\begin{longtable}{L{6cm}L{8.5cm}}
\toprule
\textbf{Test} & \textbf{Qué verifica} \\
\midrule
\endhead
\texttt{test\_\allowbreak export\_\allowbreak image\_\allowbreak png\_\allowbreak ..\_\allowbreak uses\_\allowbreak png\_\allowbreak writer\_\allowbreak and\_\allowbreak sets\_\allowbreak last\_\allowbreak dir} & Exportar PNG con nombre de archivo explícito dispara el alerta de éxito y actualiza el último directorio usado. \\
\texttt{test\_\allowbreak export\_\allowbreak image\_\allowbreak dialog\_\allowbreak cancel\_\allowbreak returns\_\allowbreak without\_\allowbreak error} & Cancelar el diálogo de guardado no dispara ningún alerta. \\
\texttt{test\_\allowbreak export\_\allowbreak image\_\allowbreak unsupported\_\allowbreak format\_\allowbreak calls\_\allowbreak error} & Un formato no soportado (\texttt{gif}) dispara \texttt{alerts.error}. \\
\texttt{test\_\allowbreak export\_\allowbreak table\_\allowbreak no\_\allowbreak chart\_\allowbreak triggers\_\allowbreak error} & Sin gráfico activo, exportar tabla dispara error. \\
\texttt{test\_\allowbreak export\_\allowbreak table\_\allowbreak chart\_\allowbreak with\_\allowbreak zero\_\allowbreak plots\_\allowbreak triggers\_\allowbreak error} & Gráfico sin curvas dispara error. \\
\texttt{test\_\allowbreak export\_\allowbreak table\_\allowbreak csv\_\allowbreak writes\_\allowbreak headers\_\allowbreak and\_\allowbreak rows} & El CSV escrito tiene el encabezado \texttt{P1\_\allowbreak X,P1\_\allowbreak Y,P2\_\allowbreak X,P2\_\allowbreak Y} y las filas correctas -- \textbf{se lee el archivo real}. \\
\texttt{test\_\allowbreak export\_\allowbreak table\_\allowbreak json\_\allowbreak writes\_\allowbreak records} & El JSON escrito tiene los registros y claves correctos -- \textbf{se lee el archivo real}. \\
\texttt{test\_\allowbreak export\_\allowbreak table\_\allowbreak uses\_\allowbreak dialog\_\allowbreak and\_\allowbreak respects\_\allowbreak cancel} & Cancelar el diálogo al exportar tabla no escribe nada ni dispara alertas. \\
\bottomrule
\end{longtable}
\end{center}

\textbf{Errores encontrados:}
\begin{itemize}
  \item \warnmark\ Los dos primeros tests de exportación de imagen \textbf{no verifican realmente que se
        haya usado el escritor PNG correcto}. El propio comentario del código lo admite:
        \textit{``Our Dummy writers are indistinguishable from outside... we rely on the absence of
        errors''}. Es decir, el test pasaría igual si el código exportara accidentalmente con el escritor de
        JPEG en vez del de PNG.
\end{itemize}

\textbf{Cómo mejorarlo:} dar a cada clase \texttt{Dummy*Writer} un atributo de clase distinto (por ejemplo
\texttt{FORMAT = "png"}) y, tras llamar a \texttt{export\_\allowbreak image}, inspeccionar qué clase fue realmente
instanciada (guardando la referencia en un contenedor mutable capturado por closure, o parcheando el
constructor para que registre la instancia creada). Así el test detectaría si el mapeo
formato→escritor se rompe.

% ------------------------------------------------------------------
\subsection{\texttt{test\_\allowbreak fileio\_\allowbreak abf.py} -- 2 tests}
% ------------------------------------------------------------------

\textbf{Qué prueba:} \texttt{FileIOService.load\_\allowbreak abf()}.

\textbf{Cómo funciona, paso a paso:}
\begin{enumerate}
  \item Una clase \texttt{DummyABF} simula el objeto que devuelve la librería \texttt{pyabf}: 2 canales de
        1000 muestras a 1000\,Hz, canal 0 = seno de 5\,Hz, canal 1 = coseno de 7\,Hz.
  \item \texttt{monkeypatch.setattr(fileio\_\allowbreak mod, "pyabf", ...)} reemplaza el módulo \texttt{pyabf} completo
        dentro de \texttt{core.services.fileio\_\allowbreak service} por un objeto que expone \texttt{ABF = DummyABF}.
  \item Se llama a \texttt{FileIOService().load\_\allowbreak abf(...)} de verdad, y se verifica la forma, la tasa de
        muestreo, los nombres de canal y valores puntuales conocidos (\texttt{sin(0)=0}, \texttt{cos(0)=1}).
  \item Una segunda clase \texttt{DummyABF\_\allowbreak Raises} simula un archivo corrupto (el constructor lanza
        \texttt{IOError}) y se verifica que el servicio propague la excepción.
\end{enumerate}

\begin{center}
\begin{longtable}{L{5.5cm}L{9cm}}
\toprule
\textbf{Test} & \textbf{Qué verifica} \\
\midrule
\endhead
\texttt{test\_\allowbreak load\_\allowbreak abf\_\allowbreak valid} & Carga correcta: forma \texttt{(2,1000)}, \texttt{fs=1000.0}, nombres de canal, valores en \texttt{t=0} correctos, eje de tiempo del mismo largo que la señal. \\
\texttt{test\_\allowbreak load\_\allowbreak abf\_\allowbreak corrupted} & Un ABF que lanza \texttt{IOError} al abrirse propaga \texttt{IOError} o \texttt{ValueError}, no lo esconde. \\
\bottomrule
\end{longtable}
\end{center}

\textbf{Errores encontrados:} ninguno funcional. Es un archivo pequeño pero completo para lo que cubre.

\textbf{Cómo mejorarlo:} existe una clase \texttt{DummyABF\_\allowbreak BadShape} (canales de largo distinto) definida
en el archivo pero \textbf{sin ningún test que la use} -- el caso ``canales con longitudes inconsistentes''
está escrito pero nunca se ejecuta. Agregar un tercer test que la use y documente si el loader debe recortar,
rellenar o lanzar una excepción en ese caso.

% ------------------------------------------------------------------
\subsection{\texttt{test\_\allowbreak fileio\_\allowbreak edf.py} -- 2 tests}
% ------------------------------------------------------------------

\textbf{Qué prueba:} \texttt{FileIOService.load\_\allowbreak edf()}. Mismo patrón que el archivo ABF, pero simulando
\texttt{pyedflib.EdfReader}.

\textbf{Cómo funciona, paso a paso:}
\begin{enumerate}
  \item \texttt{DummyEDF\_\allowbreak Uniform} simula un EDF válido de 2 canales, misma longitud (1000 muestras) y misma
        frecuencia de muestreo (1000\,Hz) en ambos.
  \item \texttt{monkeypatch} reemplaza \texttt{pyedflib.EdfReader} por esta clase.
  \item Se verifica forma, formato, nombres de canal y unidades.
  \item \texttt{DummyEDF\_\allowbreak Raises} simula un archivo corrupto; se verifica que el servicio propague la
        excepción.
\end{enumerate}

\begin{center}
\begin{longtable}{L{5.5cm}L{9cm}}
\toprule
\textbf{Test} & \textbf{Qué verifica} \\
\midrule
\endhead
\texttt{test\_\allowbreak load\_\allowbreak edf\_\allowbreak valid\_\allowbreak uniform} & Carga correcta de un EDF con canales uniformes: forma, formato, nombres y unidades. \\
\texttt{test\_\allowbreak load\_\allowbreak edf\_\allowbreak corrupted} & Un EDF que lanza \texttt{IOError} al abrirse propaga la excepción. \\
\bottomrule
\end{longtable}
\end{center}

\textbf{Errores encontrados:}
\begin{itemize}
  \item \warnmark\ El archivo define \texttt{DummyEDF\_\allowbreak Inconsistent} -- un EDF donde los dos canales tienen
        \textbf{distinta longitud y distinta frecuencia de muestreo} (1000 vs.\ 900 muestras, 1000\,Hz vs.\
        500\,Hz), con un comentario explícito: \textit{``Depending on your implementation this should raise
        (preferred) or trigger a resampling path. Here we expect a failure.''} Pero \textbf{ningún test la
        usa}. El caso más realista y más propenso a bugs (archivos EDF reales casi nunca tienen todos los
        canales a la misma frecuencia) es precisamente el que quedó sin cubrir, a pesar de que alguien ya
        había pensado en él y escrito el \textit{fixture}.
\end{itemize}

\textbf{Cómo mejorarlo:} agregar
\texttt{test\_\allowbreak load\_\allowbreak edf\_\allowbreak inconsistent\_\allowbreak channels\_\allowbreak raises\_\allowbreak or\_\allowbreak resamples} usando
\texttt{DummyEDF\_\allowbreak Inconsistent}, y decidir (y documentar en el propio \texttt{fileio\_\allowbreak service.py}) cuál de
los dos comportamientos es el esperado -- hoy no se sabe si el servicio maneja ese caso correctamente porque
nunca se ejecutó contra él.

% ------------------------------------------------------------------
\subsection{\texttt{test\_\allowbreak measurement.py} -- 28 tests}
% ------------------------------------------------------------------

\textbf{Qué prueba:} \texttt{MeasurementService} (\texttt{core/services/measurement\_\allowbreak service.py}), la
máquina de estados que maneja las mediciones interactivas de pendiente/amplitud sobre los gráficos VTK
(click en punto A, click en punto B, guardar medición).

\textbf{Cómo funciona, paso a paso:}
\begin{enumerate}
  \item \texttt{create\_\allowbreak basic\_\allowbreak service()} construye un \texttt{MeasurementService} con \texttt{parent},
        \texttt{vtk\_\allowbreak widget}, \texttt{get\_\allowbreak active\_\allowbreak chart}, \texttt{ds\_\allowbreak get} y \texttt{ds\_\allowbreak set} todos
        \texttt{Mock()} -- nunca toca VTK ni Qt real.
  \item \texttt{create\_\allowbreak service\_\allowbreak with\_\allowbreak data()} además precarga \texttt{svc.\_\allowbreak ref\_\allowbreak data} con un par de listas
        \texttt{xs}/\texttt{ys} conocidas, para los tests de búsqueda de punto más cercano.
  \item Cada grupo de tests ejercita una porción de la máquina de estados y luego \textbf{inspecciona
        atributos internos} (\texttt{svc.\_\allowbreak state}, \texttt{svc.\_\allowbreak current}, \texttt{svc.\_\allowbreak ref\_\allowbreak axes}, etc.)
        directamente -- es una prueba de caja blanca, no solo de la interfaz pública.
  \item Para las funciones que calculan la medición en sí (\texttt{two\_\allowbreak point\_\allowbreak metrics},
        \texttt{amplitude\_\allowbreak in\_\allowbreak window}), se usa \texttt{@patch(...)} para inyectar un resultado conocido y
        así aislar la lógica de guardado/estado de la lógica numérica (que vive en
        \texttt{core/filters/measurements.py} y \textbf{no tiene test propio}, ver
        sección~\ref{sec:cobertura}).
\end{enumerate}

Los 28 tests se agrupan en 8 bloques temáticos dentro del propio archivo:

\begin{center}
\begin{longtable}{L{4.5cm}p{1cm}L{7.5cm}}
\toprule
\textbf{Bloque} & \textbf{\#} & \textbf{Qué cubre} \\
\midrule
\endhead
Máquina de estados (\texttt{start}/\texttt{cancel}) & 7 & Iniciar medición de pendiente/amplitud, bloqueo por reentrada, rechazo si se pide \texttt{slope\_\allowbreak all\_\allowbreak trials} desde un plugin que no es \texttt{trials}, cancelar limpia todo el estado. \\
Persistencia (\texttt{\_\allowbreak save\_\allowbreak measurement}) & 5 & IDs autoincrementales (\texttt{slope-001, -002...}, incluso pasando de 999 a 1000), contexto (\texttt{view\_\allowbreak id}, \texttt{trial\_\allowbreak id}, \texttt{channel\_\allowbreak name}) guardado y sobreescribible. \\
Pendientes por trial & 3 & \texttt{\_\allowbreak compute\_\allowbreak and\_\allowbreak save\_\allowbreak slopes\_\allowbreak each\_\allowbreak trial\_\allowbreak by\_\allowbreak index} recorre todas las columnas \texttt{trial\_\allowbreak N} de la tabla VTK activa, calcula la pendiente de cada una y genera overlays solo para el trial actualmente mostrado. \\
Eliminación & 3 & Quitar la última medición, manejar lista vacía, limpiar todas las mediciones. \\
Helpers de búsqueda & 4 & \texttt{\_\allowbreak nearest\_\allowbreak index} (coincidencia exacta, interpolado, sin datos, con duplicados). \\
Contexto/limpieza & 2 & \texttt{cancel()} limpia todo el estado interno; \texttt{clear\_\allowbreak all\_\allowbreak measurements()} preserva el contexto (\texttt{view\_\allowbreak id}, \texttt{trial\_\allowbreak id}) aunque borre las mediciones. \\
Flujos de integración & 3 & Flujo completo de una medición de pendiente y una de amplitud de principio a fin; múltiples mediciones consecutivas incrementan el contador correctamente. \\
\bottomrule
\end{longtable}
\end{center}

\textbf{Errores encontrados:} ninguno funcional en este archivo -- es el más grande y, junto con
\texttt{test\_\allowbreak trials.py}, el más exhaustivo de la suite (28 tests para un solo servicio).

\textbf{Cómo mejorarlo:} al ser pruebas de caja blanca (dependen de nombres de atributos privados como
\texttt{\_\allowbreak state}, \texttt{\_\allowbreak current}), son frágiles ante refactors internos de \texttt{MeasurementService}
aunque el comportamiento externo no cambie. No es un error, pero conviene documentarlo como una decisión de
diseño consciente: si se refactoriza el servicio, hay que actualizar estos 28 tests aunque la interfaz
pública no cambie.

% ------------------------------------------------------------------
\subsection{\texttt{test\_\allowbreak trials.py} -- 13 tests}
% ------------------------------------------------------------------

\textbf{Qué prueba:} \texttt{core/filters/trials.py}, el módulo que detecta estímulos y corta la señal
continua en ensayos (\textit{trials}). Es el \textbf{único archivo de \texttt{services\_\allowbreak test} que no usa
ningún mock}: corre la lógica real sobre señales sintéticas (senos e impulsos) construidas a mano y sobre
objetos \texttt{SignalDataset} reales.

\textbf{Cómo funciona, paso a paso:}
\begin{enumerate}
  \item Dos generadores, \texttt{impulse\_\allowbreak train(...)} (tren de impulsos en índices dados) y
        \texttt{sine(...)} (seno de frecuencia conocida), producen señales sintéticas donde el resultado
        esperado se puede calcular a mano.
  \item \texttt{make\_\allowbreak signal\_\allowbreak dataset(...)} las empaqueta en un \texttt{SignalDataset} real.
  \item Los tests atacan el módulo en tres niveles: funciones internas de bajo nivel
        (\texttt{\_\allowbreak detect\_\allowbreak onsets\_\allowbreak abs}, \texttt{\_\allowbreak compute\_\allowbreak windows}, \texttt{\_\allowbreak extract\_\allowbreak trials\_\allowbreak matrix},
        \texttt{\_\allowbreak assemble\_\allowbreak trials\_\allowbreak general}) y la función pública de extremo a extremo
        (\texttt{cut\_\allowbreak trials\_\allowbreak single\_\allowbreak channel}).
\end{enumerate}

\begin{center}
\begin{longtable}{L{6.3cm}L{8.2cm}}
\toprule
\textbf{Test} & \textbf{Qué verifica} \\
\midrule
\endhead
\texttt{test\_\allowbreak detect\_\allowbreak onsets\_\allowbreak abs\_\allowbreak no\_\allowbreak crossings} & Señal en cero no detecta ningún cruce de umbral. \\
\texttt{test\_\allowbreak detect\_\allowbreak onsets\_\allowbreak abs\_\allowbreak single\_\allowbreak crossing} & Un único impulso se detecta en el índice correcto. \\
\texttt{test\_\allowbreak detect\_\allowbreak onsets\_\allowbreak abs\_\allowbreak multiple\_\allowbreak crossings\_\allowbreak mixed\_\allowbreak sign} & Detecta cruces positivos y negativos por igual (usa valor absoluto). \\
\texttt{test\_\allowbreak compute\_\allowbreak windows\_\allowbreak fixed\_\allowbreak basic\_\allowbreak clipping\_\allowbreak and\_\allowbreak ns} & Modo \texttt{fixed}: la ventana se recorta correctamente contra los bordes de la señal y \texttt{Ns} coincide con \texttt{(t1-t0)*fs}. \\
\texttt{test\_\allowbreak compute\_\allowbreak windows\_\allowbreak until\_\allowbreak next\_\allowbreak onset\_\allowbreak groups\_\allowbreak and\_\allowbreak last\_\allowbreak to\_\allowbreak end} & Modo \texttt{until\_\allowbreak next\_\allowbreak onset}: cada ventana termina donde empieza la siguiente, y la última llega hasta el final de la señal. \\
\texttt{test\_\allowbreak extract\_\allowbreak trials\_\allowbreak matrix\_\allowbreak alignment\_\allowbreak and\_\allowbreak padding} & El relleno (\textit{padding}) con NaN al inicio/final de cada columna respeta el \texttt{lead\_\allowbreak pad} calculado. \\
\texttt{test\_\allowbreak assemble\_\allowbreak trials\_\allowbreak general\_\allowbreak groups\_\allowbreak three\_\allowbreak per\_\allowbreak trial} & Con 3 estímulos por ensayo, agrupa correctamente 6 columnas en 2 trials. \\
\texttt{test\_\allowbreak cut\_\allowbreak trials\_\allowbreak no\_\allowbreak stimuli\_\allowbreak returns\_\allowbreak single\_\allowbreak column} & Sin estímulos detectados, devuelve 1 sola columna (toda la señal como un trial). \\
\texttt{test\_\allowbreak cut\_\allowbreak trials\_\allowbreak forced\_\allowbreak no\_\allowbreak stimulus\_\allowbreak even\_\allowbreak with\_\allowbreak triggers} & Con \texttt{stim\_\allowbreak expected=0}, ignora los disparos presentes en el canal de estímulo aunque existan. \\
\texttt{test\_\allowbreak cut\_\allowbreak trials\_\allowbreak single\_\allowbreak stim\_\allowbreak fixed\_\allowbreak window\_\allowbreak alignment\_\allowbreak and\_\allowbreak ns} & Con 2 estímulos, ventana fija: el valor de la señal en el instante del estímulo (índice cero relativo) coincide exactamente con la señal original. \\
\texttt{test\_\allowbreak cut\_\allowbreak trials\_\allowbreak single\_\allowbreak stim\_\allowbreak until\_\allowbreak next\_\allowbreak onset\_\allowbreak variable\_\allowbreak lengths} & 3 estímulos separados de forma desigual producen columnas de largo variable, rellenadas con NaN donde corresponde. \\
\texttt{test\_\allowbreak cut\_\allowbreak trials\_\allowbreak multi\_\allowbreak stim\_\allowbreak assemble\_\allowbreak ok} & 2 trials de 3 estímulos cada uno se agrupan correctamente y \texttt{stim\_\allowbreak per\_\allowbreak trial}/\texttt{stim\_\allowbreak detected} quedan bien en los metadatos. \\
\texttt{test\_\allowbreak cut\_\allowbreak trials\_\allowbreak stim\_\allowbreak channel\_\allowbreak differs\_\allowbreak from\_\allowbreak target\_\allowbreak channel} & El canal de detección de estímulo puede ser distinto del canal que se corta, y la alineación sigue siendo exacta. \\
\bottomrule
\end{longtable}
\end{center}

\textbf{Errores encontrados:} ninguno. Es, con \texttt{test\_\allowbreak dataset.py}, el ejemplo a seguir del resto de la
suite: sin mocks, con valores esperados calculados de forma independiente (no copiados de la salida del
propio código), y con nombres de test que documentan exactamente la condición de borde cubierta.

\newpage
% ==================================================================
\section{\texttt{test/plugins\_\allowbreak test/} -- validación contra referencia de MATLAB}
\label{sec:plugins-test}
% ==================================================================

Los 8 archivos de esta carpeta comparten \textbf{exactamente la misma metodología} y los mismos parámetros
de corte de ensayos -- solo cambian el plugin bajo prueba y las tolerancias numéricas.

\subsection{Metodología común, paso a paso}

\begin{enumerate}
  \item \textbf{Carga real:} \texttt{FileIOService().load\_\allowbreak abf("test/data/17308005.abf")} -- el mismo
        archivo ABF real en los 8 archivos.
  \item \textbf{Corte de trials real:} \texttt{core.filters.trials.cut\_\allowbreak trials\_\allowbreak single\_\allowbreak channel(...)} con los
        mismos parámetros en los 8 archivos: canal objetivo \texttt{0}, canal de estímulo \texttt{1}, umbral
        \texttt{0.7}, ventana \texttt{t0=-0.05\,s} a \texttt{t1=4.00\,s}, modo \texttt{until\_\allowbreak next\_\allowbreak onset}.
  \item \textbf{Doble ``sin cabeza'':} se define una subclase del plugin real (p.\ ej.\
        \texttt{AverageDouble(Average\_\allowbreak plugin)}) que sobreescribe \emph{únicamente} los métodos de
        creación de VTK/renderizado (\texttt{ensure\_\allowbreak vtk}, \texttt{render\_\allowbreak *}) para que no hagan nada, y
        expone un método \texttt{run()} que reproduce la lógica de cálculo real del botón ``Calcular'' sin
        pasar por la interfaz gráfica.
  \item \textbf{Cálculo real:} se invoca el método de cálculo real del plugin (p.\ ej.\
        \texttt{\_\allowbreak compute\_\allowbreak fft}, \texttt{compute\_\allowbreak wavelet}) -- ninguna fórmula matemática está mockeada.
  \item \textbf{Referencia MATLAB:} se carga un CSV (\texttt{test/data/*\_\allowbreak data\_\allowbreak matlab.csv}) generado
        previamente en MATLAB con el mismo ABF y los mismos parámetros.
  \item \textbf{Tres comparaciones} entre el resultado de la app y la referencia:
  \begin{enumerate}[label=(\alph*)]
    \item \textbf{Forma:} ambas matrices son 2D y no están vacías tras recortarlas al tamaño mínimo común.
    \item \textbf{Correlación de Pearson} por columna (un trial/canal) o por fila (una frecuencia, en los
          plugins de tiempo-frecuencia), exigiendo un mínimo \texttt{MIN\_\allowbreak CORR}.
    \item \textbf{Tolerancia puntual:} al menos \texttt{FRAC\_\allowbreak MIN\_\allowbreak PASS} (95\% en los 8 archivos) de los
          puntos deben cumplir \texttt{|APP - MATLAB| <= FRAC\_\allowbreak ATOL}, con un reporte detallado de las peores
          filas/columnas si falla.
  \end{enumerate}
\end{enumerate}

\subsection{Tabla de tolerancias por plugin}

\begin{center}
\begin{tabular}{lccl}
\toprule
\textbf{Plugin} & \textbf{Correlación mín.} & \textbf{Tolerancia puntual} & \textbf{Eje de comparación} \\
\midrule
\texttt{average}         & 0.98 & 0.001  & columna (trial) \\
\texttt{erp}              & 0.98 & 0.001  & columna \\
\texttt{fft}              & 0.98 & 0.75   & columna \\
\texttt{fft\_\allowbreak average}     & 0.98 & 0.15   & columna \\
\texttt{psd}              & 0.88 & 0.0001 & columna \\
\texttt{psd\_\allowbreak average}     & 0.95 & 0.0001 & columna \\
\texttt{wavelet}          & 0.80 & 0.0001 & fila (frecuencia) \\
\texttt{wavelet\_\allowbreak average} & 0.80 & 0.0001 & fila (frecuencia) \\
\bottomrule
\end{tabular}
\end{center}

Estas tolerancias no siguen ningún criterio común: \texttt{fft} acepta un error puntual de \texttt{0.75}
mientras \texttt{psd} exige \texttt{0.0001} -- cuatro órdenes de magnitud de diferencia entre dos plugins que
calculan cosas relacionadas (espectro de frecuencia). Todo indica que se ajustaron caso por caso hasta que el
test pasara, no a partir de un presupuesto de error definido de antemano.

\subsection{Error crítico \#1: los 8 archivos no se pueden ni importar}
\label{sec:import-error}

\errmark\ Los 8 archivos contienen, en sus primeras líneas, estos dos imports:

\begin{verbatim}
from core.services.fileio import FileIOService
from core.services.trial_dataset import TrialDataset
\end{verbatim}

Ninguno de esos dos módulos existe hoy en el repositorio. Los módulos reales son:

\begin{verbatim}
from core.services.fileio_service import FileIOService   # antes: core.services.fileio
from core.model.trial_dataset import TrialDataset         # antes: core.services.trial_dataset
\end{verbatim}

Al ejecutar \texttt{pytest} sobre el repositorio real, esto es lo que ocurre (salida real, verificada):

\begin{verbatim}
ERROR test/plugins_test/test_average_plugin.py
ERROR test/plugins_test/test_erp_plugin.py
ERROR test/plugins_test/test_fft_average_plugin.py
ERROR test/plugins_test/test_fft_plugin.py
ERROR test/plugins_test/test_psd_average_plugin.py
ERROR test/plugins_test/test_psd_plugin.py
ERROR test/plugins_test/test_wavelet_average_plugin.py
ERROR test/plugins_test/test_wavelet_plugin.py
8 errors during collection
\end{verbatim}

\textbf{Por qué pasó:} en algún momento se movió/renombró \texttt{core/services/fileio.py} a
\texttt{core/services/fileio\_\allowbreak service.py}, y se separó \texttt{TrialDataset} de \texttt{core/services/} a un
módulo nuevo \texttt{core/model/trial\_\allowbreak dataset.py}. Los 8 archivos de test nunca se actualizaron para
reflejar ese refactor. Como no hay integración continua (\textit{CI}) configurada en el repositorio, nadie lo
notó: la validación contra MATLAB de 8 de los plugins más importantes de la aplicación lleva rota desde ese
refactor sin que ninguna corrida automática lo haya señalado.

\textbf{Cómo mejorarlo:}
\begin{enumerate}
  \item Corregir los dos imports en los 8 archivos (cambio mecánico, una línea por archivo).
  \item Agregar un \textit{workflow} de CI (por ejemplo GitHub Actions) que corra \texttt{pytest} en cada
        \textit{push}/\textit{pull request}, para que un error de este tipo se detecte en minutos y no de
        forma manual meses después.
\end{enumerate}

\subsection{Qué se descubre al corregir el import (diagnóstico, no aplicado al repositorio)}
\label{sec:tras-fix}

Para saber si, más allá del import, el resto de la lógica de estos 8 archivos funciona, se corrigieron esos
dos imports \textbf{en una copia del código fuera del repositorio} (usada solo para diagnóstico y luego
eliminada; el repositorio real de Gamma Lab nunca fue modificado). Resultado real de esa ejecución:

\begin{center}
\begin{tabular}{lccl}
\toprule
\textbf{Archivo} & \textbf{Tests} & \textbf{Resultado} & \textbf{Detalle} \\
\midrule
\texttt{test\_\allowbreak average\_\allowbreak plugin.py} & 3 & \okmark & todos pasan \\
\texttt{test\_\allowbreak erp\_\allowbreak plugin.py} & 3 & \okmark & todos pasan \\
\texttt{test\_\allowbreak fft\_\allowbreak plugin.py} & 3 & \okmark & todos pasan \\
\texttt{test\_\allowbreak fft\_\allowbreak average\_\allowbreak plugin.py} & 3 & \okmark & todos pasan \\
\texttt{test\_\allowbreak psd\_\allowbreak plugin.py} & 3 & \okmark & todos pasan \\
\texttt{test\_\allowbreak psd\_\allowbreak average\_\allowbreak plugin.py} & 3 & \warnmark & 1 falla (correlación columna~0 = 0.58, se exige 0.95) \\
\texttt{test\_\allowbreak wavelet\_\allowbreak plugin.py} & 3 & \errmark & 2 fallan (ver \S\ref{sec:wavelet-root-cause}) \\
\texttt{test\_\allowbreak wavelet\_\allowbreak average\_\allowbreak plugin.py} & 3 & \errmark & 2 fallan (mismo patrón) \\
\midrule
\textbf{Total} & \textbf{24} & \textbf{19 pasan / 5 fallan} & \\
\bottomrule
\end{tabular}
\end{center}

\subsubsection{Causa raíz del fallo en \texttt{wavelet} y \texttt{wavelet\_\allowbreak average}}
\label{sec:wavelet-root-cause}

Los tests \texttt{test\_\allowbreak correlation\_\allowbreak by\_\allowbreak row} y \texttt{test\_\allowbreak pointwise\_\allowbreak fraction\_\allowbreak allowing\_\allowbreak outliers} de
ambos archivos fallan de forma severa: correlación mínima por fila de \textbf{-0.23} (se exige 0.80) y solo
el \textbf{0.26\%} de los puntos dentro de tolerancia (se exige 95\%). Investigando el porqué:

\begin{enumerate}
  \item Los tests usan por defecto \texttt{scale\_\allowbreak log=True} (remuestreo del eje de frecuencia a escala
        logarítmica antes de comparar), igual que la interfaz cuando el usuario activa la casilla ``Log''.
  \item Al desactivar únicamente ese paso (\texttt{\_\allowbreak scale\_\allowbreak log}, en
        \texttt{wavelet\_\allowbreak plugin.py}, líneas 242--267) y comparar el escalograma \emph{crudo} contra la misma
        referencia de MATLAB, la correlación media por fila sube de \textbf{0.22 a 0.75}. Esto aísla a
        \texttt{\_\allowbreak scale\_\allowbreak log} como el responsable principal del colapso: el remuestreo interpola
        linealmente sobre una rejilla de solo 998 puntos (\textasciitilde0.5\,Hz de paso), lo cual degrada
        gravemente la resolución cerca de 1\,Hz.
  \item Aun así, con el escalograma crudo (sin ese paso), el \textbf{60\% de las filas (596 de 998)} siguen
        por debajo del umbral de correlación, y la razón de magnitud APP/MATLAB varía fila a fila entre
        \textbf{0.54$\times$ y 11.3$\times$} -- no es un factor de escala constante corregible con un
        multiplicador único. Esto es compatible con dos causas adicionales ya identificadas en el código:
        la decimación de la señal (\texttt{sig[::factor]}, línea 202) \textbf{sin filtro anti-aliasing}
        previo, y una posible diferencia de convención de normalización entre el wavelet Morlet de
        \texttt{pywt.cwt} y el \texttt{cwt} de MATLAB.
\end{enumerate}

\textbf{Cómo mejorarlo:}
\begin{enumerate}
  \item Revisar y corregir \texttt{\_\allowbreak scale\_\allowbreak log} en \texttt{wavelet\_\allowbreak plugin.py} -- probablemente necesita una
        rejilla de interpolación más fina que la resolución lineal original, no la misma cantidad de puntos.
  \item Agregar un filtro pasa-bajos antes de cualquier decimación en \texttt{compute\_\allowbreak wavelet}.
  \item Verificar contra la documentación de MATLAB qué convención de normalización usa su \texttt{cwt} para
        el wavelet de Morlet, y ajustar el cálculo de \texttt{pywt} para que coincida (factor de escala
        explícito si es necesario).
  \item Una vez corregido, bajar \texttt{FRAC\_\allowbreak ATOL} y subir \texttt{MIN\_\allowbreak CORR} a valores acordes al resto de
        la suite (hoy estos dos plugins tienen, con diferencia, las tolerancias más laxas) para que el test
        siga siendo una guardia de regresión útil y no solo un test que ``pasa como sea''.
\end{enumerate}

\subsubsection{Fallo en \texttt{psd\_\allowbreak average}}

La columna~0 (el primer trial) tiene correlación 0.58 contra un umbral de 0.95. Es consistente con un
hallazgo de corrección independiente detectado al leer el código:
\texttt{psd\_\allowbreak average\_\allowbreak plugin.py}, línea 255, llama a
\texttt{np.nan\_\allowbreak to\_\allowbreak num(Xds, copy=False, ...)} -- cuando no hay decimación (\texttt{srt==1}),
\texttt{Xds} es literalmente el mismo array en memoria que \texttt{X}, así que esta limpieza de NaN
\textbf{muta el array de trials compartido en el sitio}, en vez de trabajar sobre una copia. Cualquier otro
plugin que use el mismo \texttt{TrialDataset} después puede estar viendo datos ya alterados.

\textbf{Cómo mejorarlo:} cambiar la línea 255 a
\texttt{Xds = np.nan\_\allowbreak to\_\allowbreak num(Xds, copy=True, ...)} (o copiar \texttt{X} antes de decimar), y agregar un test
unitario en \texttt{services\_\allowbreak test} que verifique explícitamente que \texttt{\_\allowbreak compute\_\allowbreak psd} no modifica su
argumento de entrada.

\newpage
% ==================================================================
\section{Detalle por archivo de \texttt{test/plugins\_\allowbreak test/}}
% ==================================================================

Todos siguen la plantilla de la sección anterior; aquí se detalla qué calcula específicamente cada plugin y
sus 3 tests.

\subsection{\texttt{test\_\allowbreak average\_\allowbreak plugin.py}}
Compara el promedio punto a punto entre trials (\texttt{np.mean} sobre el eje de trials) contra
\texttt{average\_\allowbreak data\_\allowbreak matlab.csv}.
\begin{center}
\begin{tabular}{ll}
\toprule
\texttt{test\_\allowbreak shape\_\allowbreak and\_\allowbreak length} & Forma no vacía tras recorte a tamaño común. \\
\texttt{test\_\allowbreak correlation\_\allowbreak with\_\allowbreak matlab} & Correlación $\geq 0.98$. \\
\texttt{test\_\allowbreak pointwise\_\allowbreak fraction\_\allowbreak allowing\_\allowbreak outliers} & $\geq$95\% de puntos con error $\leq 0.001$. \\
\bottomrule
\end{tabular}
\end{center}
Estado real: \okmark\ los 3 pasan (una vez corregido el import de la sección~\ref{sec:import-error}).

\subsection{\texttt{test\_\allowbreak erp\_\allowbreak plugin.py}}
Compara la matriz \textit{butterfly} (una fila por trial) contra \texttt{erp\_\allowbreak data\_\allowbreak matlab.csv}.
\begin{center}
\begin{tabular}{ll}
\toprule
\texttt{test\_\allowbreak shape\_\allowbreak and\_\allowbreak columns} & Forma no vacía tras recorte. \\
\texttt{test\_\allowbreak correlation\_\allowbreak by\_\allowbreak column} & Correlación por trial $\geq 0.98$. \\
\texttt{test\_\allowbreak pointwise\_\allowbreak fraction\_\allowbreak allowing\_\allowbreak outliers} & $\geq$95\% de puntos con error $\leq 0.001$. \\
\bottomrule
\end{tabular}
\end{center}
Estado real: \okmark\ los 3 pasan.

\subsection{\texttt{test\_\allowbreak fft\_\allowbreak plugin.py}}
Compara la magnitud de \texttt{np.fft.rfft} (sin ventana, sin normalizar) contra \texttt{fft\_\allowbreak data\_\allowbreak matlab.csv}.
\begin{center}
\begin{tabular}{ll}
\toprule
\texttt{test\_\allowbreak shape\_\allowbreak and\_\allowbreak columns} & Forma no vacía tras recorte. \\
\texttt{test\_\allowbreak correlation\_\allowbreak by\_\allowbreak column} & Correlación por trial $\geq 0.98$. \\
\texttt{test\_\allowbreak pointwise\_\allowbreak fraction\_\allowbreak allowing\_\allowbreak outliers} & $\geq$95\% de puntos con error $\leq 0.75$ (tolerancia más laxa de toda la suite). \\
\bottomrule
\end{tabular}
\end{center}
Estado real: \okmark\ los 3 pasan -- pero con una tolerancia puntual de \texttt{0.75} que oculta casi
cualquier discrepancia salvo las más groseras.

\subsection{\texttt{test\_\allowbreak fft\_\allowbreak average\_\allowbreak plugin.py}}
Igual que el anterior, pero sobre el promedio de la magnitud FFT entre trials; referencia
\texttt{fft\_\allowbreak average\_\allowbreak data\_\allowbreak matlab.csv}. Estado real: \okmark\ los 3 pasan (tolerancia \texttt{0.15}).

\subsection{\texttt{test\_\allowbreak psd\_\allowbreak plugin.py}}
Compara la densidad espectral de potencia (\texttt{scipy.signal.welch}) contra \texttt{psd\_\allowbreak data\_\allowbreak matlab.csv}.
\begin{center}
\begin{tabular}{ll}
\toprule
\texttt{test\_\allowbreak shape\_\allowbreak and\_\allowbreak columns} & Forma no vacía tras recorte. \\
\texttt{test\_\allowbreak correlation\_\allowbreak by\_\allowbreak column} & Correlación por trial $\geq 0.88$. \\
\texttt{test\_\allowbreak pointwise\_\allowbreak fraction\_\allowbreak allowing\_\allowbreak outliers} & $\geq$95\% de puntos con error $\leq 0.0001$. \\
\bottomrule
\end{tabular}
\end{center}
Estado real: \okmark\ los 3 pasan.

\subsection{\texttt{test\_\allowbreak psd\_\allowbreak average\_\allowbreak plugin.py}}
Igual que el anterior sobre el promedio entre trials; referencia \texttt{psd\_\allowbreak average\_\allowbreak data\_\allowbreak matlab.csv}.
Estado real: \warnmark\ \texttt{test\_\allowbreak shape\_\allowbreak and\_\allowbreak columns} y \texttt{test\_\allowbreak pointwise\_\allowbreak fraction\_\allowbreak allowing\_\allowbreak outliers}
pasan; \textbf{\texttt{test\_\allowbreak correlation\_\allowbreak by\_\allowbreak column} falla} (columna~0, corr.\ 0.58 vs.\ umbral 0.95). Ver
causa raíz en la sección anterior (mutación in-place de NaN).

\subsection{\texttt{test\_\allowbreak wavelet\_\allowbreak plugin.py}}
Compara el escalograma de un único trial (CWT con wavelet de Morlet vía \texttt{pywt.cwt}) contra
\texttt{wavelet\_\allowbreak data\_\allowbreak matlab.csv}.
\begin{center}
\begin{tabular}{ll}
\toprule
\texttt{test\_\allowbreak shape\_\allowbreak and\_\allowbreak size} & Forma no vacía tras recorte. \\
\texttt{test\_\allowbreak correlation\_\allowbreak by\_\allowbreak row} & Correlación por frecuencia $\geq 0.80$. \\
\texttt{test\_\allowbreak pointwise\_\allowbreak fraction\_\allowbreak allowing\_\allowbreak outliers} & $\geq$95\% de puntos con error $\leq 0.0001$. \\
\bottomrule
\end{tabular}
\end{center}
Estado real: \errmark\ solo \texttt{test\_\allowbreak shape\_\allowbreak and\_\allowbreak size} pasa; los otros dos fallan. Ver
sección~\ref{sec:wavelet-root-cause}.

\subsection{\texttt{test\_\allowbreak wavelet\_\allowbreak average\_\allowbreak plugin.py}}
Igual que el anterior, promediando el escalograma de todos los trials; referencia
\texttt{wavelet\_\allowbreak average\_\allowbreak data\_\allowbreak matlab.csv}. Estado real: \errmark\ mismo patrón de fallo que
\texttt{wavelet}.

\newpage
% ==================================================================
\section{\texttt{test/performance\_\allowbreak test/} -- benchmarks de rendimiento}
\label{sec:performance}
% ==================================================================

Carpeta agregada durante esta auditoría (no existía antes). No verifica corrección -- mide tiempo de
ejecución del código \textbf{real}, sin mocks, contra el ABF real, para poder comparar ``antes'' y
``después'' de futuras optimizaciones de rendimiento.

\subsection{Infraestructura compartida: \texttt{bench\_\allowbreak common.py}}

No es un archivo de test en sí, sino el módulo de utilidades que usan los otros dos:

\begin{enumerate}
  \item \texttt{time\_\allowbreak block(fn, *args, reps=5, warmup=1, **kwargs)} -- ejecuta \texttt{fn} una vez de
        calentamiento (se descarta) y luego \texttt{reps} veces cronometradas con
        \texttt{time.perf\_\allowbreak counter()}; devuelve mínimo, mediana, media y máximo en segundos.
  \item \texttt{record(nombre, stats, label=...)} -- agrega el resultado a
        \texttt{results/perf\_\allowbreak results.json}, bajo la etiqueta del \textit{run} (variable de entorno
        \texttt{GAMMA\_\allowbreak PERF\_\allowbreak LABEL}, por defecto \texttt{"baseline"}). Etiquetas distintas no se pisan entre
        sí.
  \item \texttt{compare\_\allowbreak report.py} -- script standalone que lee ese JSON y imprime una tabla
        antes/después con el factor de aceleración (\textit{speedup}) por benchmark.
\end{enumerate}

\subsection{\texttt{test\_\allowbreak compute\_\allowbreak bench.py} -- 7 tests}

\textbf{Qué mide:} el tiempo de cómputo puro (sin VTK) de los 6 plugins de análisis. Hoy, 5 de esos 6
(\texttt{fft}, \texttt{fft\_\allowbreak average}, \texttt{psd}, \texttt{psd\_\allowbreak average}, \texttt{relative\_\allowbreak psd}) y el
\texttt{wavelet} simple corren \textbf{en el hilo principal de Qt}, sin \texttt{QThread} -- así que este
número es directamente ``cuánto se congela la interfaz por cada clic en Calcular''.

\textbf{Cómo funciona, paso a paso:}
\begin{enumerate}
  \item \textit{Fixtures} de módulo cargan el ABF real una sola vez y cortan los trials con los mismos
        parámetros que \texttt{plugins\_\allowbreak test} (sección~\ref{sec:plugins-test}).
  \item Cada test instancia el plugin real (sin interfaz gráfica) y llama a \texttt{time\_\allowbreak block} sobre su
        método de cálculo real (p.\ ej.\ \texttt{plug.\_\allowbreak compute\_\allowbreak fft}, \texttt{plug.compute\_\allowbreak wavelet}).
  \item Los resultados se imprimen en vivo con \texttt{pytest -s} y quedan guardados en el JSON.
\end{enumerate}

\begin{center}
\begin{tabular}{lr}
\toprule
\textbf{Benchmark} & \textbf{Mediana medida} \\
\midrule
\texttt{fft.\_\allowbreak compute\_\allowbreak fft} (60 trials) & 3.7\,ms \\
\texttt{fft\_\allowbreak average.\_\allowbreak compute\_\allowbreak fft\_\allowbreak average} (60 trials) & 8.3\,ms \\
\texttt{psd.\_\allowbreak compute\_\allowbreak psd} (Welch, 60 trials) & 4.9\,ms \\
\texttt{psd\_\allowbreak average.\_\allowbreak compute\_\allowbreak psd} (Welch, 60 trials) & 9.6\,ms \\
\texttt{relative\_\allowbreak psd.\_\allowbreak compute\_\allowbreak psd} (Welch, 60 trials) & 5.7\,ms \\
\texttt{relative\_\allowbreak psd.\_\allowbreak compute\_\allowbreak relative\_\allowbreak psd} (razón) & 0.07\,ms \\
\texttt{wavelet.compute\_\allowbreak wavelet} (1 trial, CWT) & \textbf{203\,ms} \\
\texttt{wavelet\_\allowbreak average.compute\_\allowbreak wavelet} ($\times$20 trials) & \textbf{4.17\,s} \\
\bottomrule
\end{tabular}
\end{center}

\textbf{``Error'' encontrado (de firma, no de lógica):} al escribirlo, \texttt{psd\_\allowbreak average.\_\allowbreak compute\_\allowbreak psd}
resultó tener \textbf{una firma distinta} a \texttt{psd}/\texttt{relative\_\allowbreak psd} (no acepta el parámetro
\texttt{detrend}) -- confirmado al leer el código, no es un bug del test sino una inconsistencia real entre
plugins que deberían ser gemelos. El benchmark se ajustó para llamarlo con su firma real.

\textbf{Cómo mejorarlo (del código bajo prueba, no del test):} unificar la firma de \texttt{\_\allowbreak compute\_\allowbreak psd}
entre \texttt{psd}, \texttt{psd\_\allowbreak average} y \texttt{relative\_\allowbreak psd} -- hoy son tres copias casi idénticas con
pequeñas diferencias accidentales, lo cual además facilita que un fix se aplique a una copia y se olvide en
las otras (exactamente lo que pasó con el bug de NaN de la sección~\ref{sec:tras-fix}).

\subsection{\texttt{test\_\allowbreak render\_\allowbreak bench.py} -- 3 tests}

\textbf{Qué mide:} el bucle \texttt{for i... for j...: img.SetScalarComponentFromFloat(...)} que construyen
\texttt{wavelet}, \texttt{wavelet\_\allowbreak average} y \texttt{erp} para dibujar su mapa 2D en VTK, en vez de
vectorizar con \texttt{vtk.util.numpy\_\allowbreak support.numpy\_\allowbreak to\_\allowbreak vtk} (como ya hace
\texttt{core/utils/adapters.py} para todos los gráficos de líneas de la app).

\textbf{Cómo funciona, paso a paso:}
\begin{enumerate}
  \item Se instancia el plugin real y se le asignan los atributos mínimos que \texttt{render\_\allowbreak scalogram}/
        \texttt{\_\allowbreak render\_\allowbreak heatmap} necesitan (\texttt{vtk\_\allowbreak widget}, \texttt{\_\allowbreak context\_\allowbreak view}) apuntando a
        una \texttt{vtkRenderWindow} en modo \textit{offscreen} (\texttt{SetOffScreenRendering(1)}) -- no se
        abre ninguna ventana, pero el código de renderizado que corre es exactamente el real.
  \item Se cronometra la llamada real a \texttt{render\_\allowbreak scalogram}/\texttt{\_\allowbreak render\_\allowbreak heatmap} con datos
        reales (un escalograma calculado con \texttt{compute\_\allowbreak wavelet} sobre el ABF real).
\end{enumerate}

\begin{center}
\begin{tabular}{lrr}
\toprule
\textbf{Benchmark} & \textbf{Tamaño} & \textbf{Mediana} \\
\midrule
\texttt{wavelet.render\_\allowbreak scalogram} & 998$\times$3051 (3.0M puntos) & \textbf{701\,ms} \\
\texttt{wavelet\_\allowbreak average.render\_\allowbreak scalogram} & 998$\times$3051 (3.0M puntos) & \textbf{728\,ms} \\
\texttt{erp.\_\allowbreak render\_\allowbreak heatmap} & 60$\times$30501 $\to$ recortado internamente a 60$\times$1907 & 34\,ms \\
\bottomrule
\end{tabular}
\end{center}

\textbf{Por qué \texttt{erp} sale tan barato en comparación:} \texttt{erp\_\allowbreak plugin.py} tiene un límite interno
(\texttt{MAX\_\allowbreak SAMPLES = 2000}) que recorta el eje de tiempo \emph{antes} de entrar al bucle, así que en la
práctica solo itera sobre $\sim$114\,000 puntos, no 1.8 millones. Sufre el mismo defecto de diseño que
\texttt{wavelet}/\texttt{wavelet\_\allowbreak average}, pero un límite de resolución ya existente lo protege de pagar el
costo completo.

\textbf{Cómo mejorarlo:} reemplazar el doble bucle de Python en los tres archivos por
\texttt{numpy\_\allowbreak support.numpy\_\allowbreak to\_\allowbreak vtk(Z.ravel(order="C"))} seguido de
\texttt{img.GetPointData().SetScalars(vtk\_\allowbreak array)} -- el mismo patrón que ya usa
\texttt{core/utils/adapters.py}. Es una función existente y probada en el propio repositorio; el problema no
es la ausencia de la herramienta correcta, sino que estos tres archivos nunca la adoptaron.

\newpage
% ==================================================================
\section{Cobertura: qué no tiene ningún test}
\label{sec:cobertura}
% ==================================================================

No hay archivo de test para:

\begin{multicols}{2}
\begin{itemize}[nosep]
  \item \texttt{core/kernel.py}
  \item \texttt{core/plugins/loader.py}
  \item \texttt{core/plugins/manager.py}
  \item \texttt{core/plugins/interfaces.py}
  \item \texttt{core/services/settings\_\allowbreak service.py}
  \item \texttt{core/utils/adapters.py}
  \item \texttt{core/utils/plugin\_\allowbreak alerts.py}
  \item \texttt{core/utils/vtk\_\allowbreak context\_\allowbreak menu.py}
  \item \texttt{core/filters/measurements.py}
  \item \texttt{app/view/main\_\allowbreak window.py}
  \item \texttt{plugins/preprocessing/prepare/filter/}
  \item \texttt{plugins/preprocessing/prepare/artifact\_\allowbreak remove/}
  \item \texttt{plugins/preprocessing/trials/} (el plugin de UI, solo el módulo interno \texttt{core.filters.trials} tiene test)
  \item \texttt{plugins/io/open\_\allowbreak signal/}
  \item \texttt{plugins/measure/amplitude/}
  \item \texttt{plugins/measure/slope/}
  \item \texttt{plugins/analysis/frequency/relative\_\allowbreak psd/} (sección \S\ref{sec:relative-psd-bug})
  \item \texttt{plugins/faq/help/}
  \item Todos los archivos \texttt{*\_\allowbreak ui.py} de cada plugin
\end{itemize}
\end{multicols}

Sobre este último punto: \texttt{relative\_\allowbreak psd} no tiene ningún archivo en \texttt{test/plugins\_\allowbreak test/} --
fue precisamente el plugin cuyo comportamiento revisamos manualmente en esta auditoría (el selector
``Calculation Mode (GF)'' no tiene ningún efecto posible sobre el resultado de Potencia Relativa, por
construcción algebraica de la fórmula en \texttt{\_\allowbreak compute\_\allowbreak relative\_\allowbreak psd}). Es un ejemplo concreto de cómo
un hueco de cobertura se traduce directamente en un bug que llegó al usuario final sin que ningún test lo
hubiera podido atrapar.
\label{sec:relative-psd-bug}

\newpage
% ==================================================================
\section{Tabla resumen de todos los errores encontrados}
% ==================================================================

\begin{center}
\begin{longtable}{L{1.6cm}L{4.3cm}L{8.3cm}}
\toprule
\textbf{Severidad} & \textbf{Ubicación} & \textbf{Descripción} \\
\midrule
\endhead
\errmark & \texttt{test/plugins\_\allowbreak test/*.py} (8 archivos) & Import roto: \texttt{core.services.fileio} / \texttt{core.services.trial\_\allowbreak dataset} no existen. 0 de 24 tests se ejecutan. \\
\errmark & \texttt{wavelet\_\allowbreak plugin.py} \texttt{\_\allowbreak scale\_\allowbreak log} & Remuestreo a escala log degrada la correlación con MATLAB de 0.75 a 0.22; bug real de implementación, no solo de tolerancia. \\
\errmark & \texttt{wavelet\_\allowbreak plugin.py} \texttt{compute\_\allowbreak wavelet} & Decimación sin filtro anti-aliasing + posible desajuste de normalización Morlet vs.\ MATLAB; 60\% de filas por debajo del umbral incluso sin el bug de \texttt{\_\allowbreak scale\_\allowbreak log}. \\
\errmark & \texttt{psd\_\allowbreak average\_\allowbreak plugin.py}:255 & \texttt{np.nan\_\allowbreak to\_\allowbreak num(Xds, copy=False)} muta en el sitio el array de trials compartido; probable causa del fallo de correlación en columna~0. \\
\warnmark & \texttt{test\_\allowbreak export.py} & Los tests de exportación PNG no verifican qué escritor VTK se usó realmente, solo que no hubo excepción. \\
\warnmark & \texttt{test\_\allowbreak fileio\_\allowbreak edf.py} & \texttt{DummyEDF\_\allowbreak Inconsistent} (canales de distinta frecuencia) está escrita pero ningún test la usa. \\
\warnmark & \texttt{test\_\allowbreak fileio\_\allowbreak abf.py} & \texttt{DummyABF\_\allowbreak BadShape} está escrita pero ningún test la usa. \\
\warnmark & Tolerancias de \texttt{plugins\_\allowbreak test} & Sin criterio común (0.0001 a 0.75); parecen ajustadas para pasar, no derivadas de un presupuesto de error. \\
\warnmark & \texttt{psd\_\allowbreak average.\_\allowbreak compute\_\allowbreak psd} & Firma distinta a \texttt{psd}/\texttt{relative\_\allowbreak psd} (sin parámetro \texttt{detrend}) sin razón aparente. \\
\warnmark & Toda la suite & Sin integración continua (CI): el import roto de 8 archivos completos no se detectó automáticamente. \\
\warnmark & \S\ref{sec:cobertura} & 12 archivos de \texttt{core/} y 8 plugins sin ningún test, incluyendo \texttt{relative\_\allowbreak psd}, donde ya se confirmó un bug real. \\
\bottomrule
\end{longtable}
\end{center}

% ==================================================================
\section{Recomendaciones priorizadas}
% ==================================================================

\begin{enumerate}
  \item \textbf{Corregir los dos imports} en los 8 archivos de \texttt{test/plugins\_\allowbreak test/} -- cambio de una
        línea por archivo, restaura de inmediato el 100\% de la validación contra MATLAB.
  \item \textbf{Agregar un \textit{workflow} de CI} que corra \texttt{pytest} en cada cambio, para que este
        tipo de rotura se detecte en minutos, no en la siguiente auditoría manual.
  \item \textbf{Investigar y corregir \texttt{\_\allowbreak scale\_\allowbreak log} y la decimación sin filtro} en
        \texttt{wavelet\_\allowbreak plugin.py}/\texttt{wavelet\_\allowbreak average\_\allowbreak plugin.py} -- es el fallo más severo
        encontrado (correlación negativa contra la referencia).
  \item \textbf{Corregir la mutación in-place de NaN} en \texttt{psd\_\allowbreak average\_\allowbreak plugin.py:255} -- copiar el
        array antes de limpiar.
  \item \textbf{Vectorizar el renderizado} de \texttt{wavelet}, \texttt{wavelet\_\allowbreak average} y \texttt{erp} con
        \texttt{numpy\_\allowbreak to\_\allowbreak vtk}, reutilizando el patrón que ya existe en \texttt{core/utils/adapters.py}
        (fix de rendimiento medido en \S\ref{sec:performance}: hasta 700\,ms por cálculo).
  \item \textbf{Cubrir con tests} el kernel, el sistema de plugins, \texttt{filter}, \texttt{artifact\_\allowbreak remove},
        \texttt{trials\_\allowbreak plugin}, \texttt{open\_\allowbreak signal}, \texttt{amplitude}/\texttt{slope} y, en particular,
        \texttt{relative\_\allowbreak psd} -- donde ya se confirmó un bug de comportamiento sin ningún test que lo
        hubiera atrapado.
  \item \textbf{Revisar las tolerancias} de \texttt{fft}/\texttt{fft\_\allowbreak average}/\texttt{wavelet} una vez
        corregidos los bugs de fondo, para que dejen de estar ajustadas ``a que pase'' y vuelvan a ser una
        guardia de regresión real.
  \item \textbf{Completar los casos de borde ya escritos pero sin usar:} \texttt{DummyEDF\_\allowbreak Inconsistent} y
        \texttt{DummyABF\_\allowbreak BadShape} solo necesitan un test que los invoque.
\end{enumerate}

% ==================================================================
\section{Conclusión}
% ==================================================================

La suite de pruebas de Gamma Lab tiene, en las partes que sí se ejecutan
(\texttt{services\_\allowbreak test} y la nueva \texttt{performance\_\allowbreak test}), un diseño sólido: datos reales o sintéticos
con resultado esperado calculado de forma independiente, mocking quirúrgico solo donde hace falta (VTK, Qt,
\texttt{pyabf}, \texttt{pyedflib}), y aserciones concretas. El problema no es de diseño sino de
\textbf{mantenimiento}: la capa que más importa para la credibilidad científica de la aplicación --la
validación de cada plugin de análisis contra una referencia de MATLAB-- lleva rota por un import desde un
refactor que nadie sincronizó con los tests, y nadie lo detectó porque no hay integración continua. Al
corregir ese único punto de falla, la suite deja de ``pasar en silencio'' y empieza a hacer su trabajo:
señala tres plugins (\texttt{wavelet}, \texttt{wavelet\_\allowbreak average}, \texttt{psd\_\allowbreak average}) con errores
numéricos reales que hoy llegan al usuario sin ninguna alerta.

\end{document}
